\documentclass[a4paper,11pt]{article}
\usepackage[utf8]{inputenc}
\usepackage[T1]{fontenc}
\usepackage{enumitem}
\usepackage{geometry}
\geometry{margin=2.5cm}

\title{Business Need Statement\\\large Management of Individual and Group Training Sessions}

\begin{document}
\maketitle

\section{Problem}
Conducting individual and group training sessions in swimming pools and other sports environments requires the trainer or session instructor to demonstrate exercises, follow the training program, and provide instructions to participants at the same time.

In a swimming pool, the instructor is usually positioned at the poolside and therefore cannot always properly demonstrate movements performed in the water or underwater. Participants may find it difficult to clearly see the correct body position, sequence of movements, or exercise technique.

Communication is also affected by noise, distance, and the number of participants. Not all participants can hear the instructor equally well, which in practice limits the size of a group that can be managed effectively at the same time.

The aim is to reduce these limitations related to visual demonstration and audio communication and to make the organization of individual and group training sessions easier.

\section{Solution}
The system should allow training programs to be prepared in advance using different stages, video content, and audio instructions. During a training session, the prepared program would run according to the predefined sequence, while the instructor could control it, adjust its progress, or move to another stage when necessary.

Participants would receive the required information through visual and audio channels. Visual content could demonstrate movements that the instructor cannot properly show from their current position, while audio instructions should be clearly audible to all participants.

The solution should be adaptable to swimming pools and other sports environments, different group sizes, and various types of individual and group training sessions.

\section{Business Benefits}

The virtual trainer mode will allow prepared training sessions to run during otherwise unused pool or sports facility time without requiring the instructor to remain on site. This will improve facility utilization, make more sessions available, and reduce the labor cost of an individual session.

If the advertising extension is implemented, screens may display advertising or information between sessions and unobtrusive sponsor logos during a session. Possible advertising revenue could help recover the investment in the system and hardware; this function is outside the baseline.

Training quality will be improved through precise demonstrations of exercises in the water and underwater, together with individual audio delivery through wireless headphones. Participants will hear instructions clearly without disturbing other pool users, and group size will no longer be constrained by the audibility of the instructor's voice.

After an individual assessment, the instructor will be able to prepare a personalized program, assign preparation activities or exercises to be completed at home, and make them available through the responsive web interface on a phone. The program may be updated remotely without an additional meeting, while the system will track bookings and actual attendance. Progress analysis is an extension.

\section{Needs Included in the Project}
\begin{enumerate}[label=\arabic*.]
\item Planning and organization of individual and group training sessions.
\item Creation of training programs consisting of different stages, exercises, and other content.
\item Presentation of visual information to participants during training sessions.
\item Delivery of audio instructions and the instructor's voice to participants.
\item Automatic execution of a prepared training program with the possibility for the instructor to control its progress.
\item Management of participants and their participation in training sessions.
\item Management of training content, programs, and their history.
\item Use of the system in swimming pools and indoor group fitness classes: in a pool, exercises are shown on a suitable screen and audio is delivered through participant headphones; in a gym, a TV on wheels may show the exercises while suitable room audio equipment delivers instructions. The session setup must support different group sizes and the video and audio equipment available at each venue.
\item Convenient viewing of the training schedule, place booking, and registration management.
\item Recording of actual attendance and participation history.
\item A virtual trainer mode that allows a prepared session to run without the instructor being continuously present on site.
\item Creation of personalized client programs, preparation recommendations, and exercises to be completed at home.
\item The ability for the instructor to update a client's program and replace exercises remotely without an additional meeting.
\end{enumerate}

\section{Desired Extensions}
The following needs are considered possible extensions of the system and are not required for the initial version:
\begin{enumerate}[label=\arabic*., start=14]
\item A mobile application for participants and instructors in which clients can view their program, preparation recommendations, and exercises to be completed at home.
\item Identification of participants in the sports environment and association of participants with a specific training session.
\item Creation of a library of different training programs and video content.
\item Adaptation of training programs according to participant needs, fitness level, or previous activity.
\item Analysis of participant progress and training results.
\item Use of the system across multiple locations or different sports centers.
\item Display of advertising and informational content on screens outside training sessions and sponsor logos during a session.
\end{enumerate}

\end{document}
